% SPDX-FileCopyrightText: 2026 Bart van de Lint
% SPDX-License-Identifier: AGPL-3.0-or-later
\section{The model against measurements}
\label{app:validation}
\label{sec:checks}

This appendix sets the model beside the strongest published measurements it can be held against: whole ships at sea first, since they test every part of the model together, then the weather the ships sail in, then the parts. Each gap is the model less the measurement, so a positive gap means the model gives more. \Cref{tab:validation} lists every check in one line; the text keeps those a reader needs in order to trust the comparison between the rigs.

\begin{table*}[tbp]
\centering
\footnotesize
\caption{\given{Each check of the model against measurement: what is compared, the measurement and its source, and the gap, the model less the measurement. Where a model value was set from the same measurement, the row says so and the gap shows only how far single units scatter about it.}}
\label{tab:validation}
\given{%
\begin{tabular}{@{}l>{\raggedright\arraybackslash}p{4.6cm}>{\raggedright\arraybackslash}p{5.4cm}>{\raggedright\arraybackslash}p{4.4cm}@{}}
\toprule
Part & Compared & Measured & Gap \\
\midrule
At sea & Rotors' fuel saving & two ships over months of service \citep{lr2025trlady,norsepower2019pelican} & $+$\NumMeasuredRotorGap{} points \\
 & Wing sails' fuel saving & two ships \citep{gcaptain2024pyxis,mol2024shofu} & $-$\NumMeasuredPyxisGap{} points, and below a range \\
 & Kite's fuel saved per day & the Beluga SkySails' first Atlantic crossing \citep{verkehrsrundschau2008beluga} & \NumBelugaTonnes{} against 0.5~t a day \\
 & Rotor thrust & speed trial of the ferry Copenhagen \citep{werner2023copenhagen} & within \NumCopenhagenThrustBow\,\% with the wind ahead, $+$\NumCopenhagenThrustBeam{} to $+$\NumCopenhagenThrustAft\,\% abeam and aft \\
\addlinespace
Weather & ERA5 wind at \SI{10}{m} & six moored buoys on the routes, 2017 to 2019 \citep{ndbc2026stdmet} & \NumEraWindBias~m/s low; \NumEraStrongWind\,\% low from \SI{12}{m/s} \\
 & ERA5 wave height & the same buoys & within \NumEraWaveBias~m; \NumEraHighWaves\,\% low from \SI{3}{m} \\
\addlinespace
Ship & Calm-water resistance & tank tests of a bulk carrier, a tanker and a container ship \citep{t2015,farkas2017} & $+0.3$ to $+10$\,\%, growing with speed \\
 & Calm-water power, uncalibrated & source designs of the classes \citep{man2019mr,man2022bulk,man2024container} & $-6$ to $+32$\,\%; one factor per class removes it \\
 & Coaster engine & Conoship CIP3800 and Damen Combi Freighter 3850 & $+11$\,\% \\
 & Added resistance in waves & \num{1127} model-test samples \citep{mittendorf2022} & inside their scatter \\
 & Hull side force at leeway & tank tests of two hulls without a skeg \citep{vanderkolk2021} & $+40$ to $+70$\,\% \\
 & Drag of the leeway & the same tests, carried to full scale \citep{kramer2022b} & about $-50$\,\% \\
 & Centre of the hull's side force & captive tests of a tanker \citep{toxopeus2011} & within the scatter of repeated runs \\
 & Rudder drag & lifting-line model against CFD \citep{kramer2022b} & within 10\,\% sailing straight, 25\,\% on average \\
 & Hull forces at heel & CFD of a cargo ship heeled \SI{20}{\degree} \citep{kramer2022b} & none modelled; side force 7 to 20\,\% lower there \\
 & Engine fuel at part load & MAN L21/31 test-bed fuel curves \citep{man2024l2131} & $+1$ to $+7$\,\% at 50 to 85\,\% load \\
 & Share of the fuel saving the side force takes & ships with rotors and sails modelled with it \citep{tillig2019,kramer2022} & \cref{sec:val_drift} \\
\addlinespace
Rigs & Wing sail highest lift & wind tunnel and CFD of a two-element wing \citep{wright2024wingsail} & $+5$ to $+11$\,\%, the flap held \\
 & Suction sail force & strain gauges on two VentiFoils at sea \citep{ghorbani2024suction}, which set it & 0 \\
 & Rotor lift & ShipCLEAN fit to the Copenhagen's trial \citep{tillig2020} & $-4$\,\% \\
 & Losses between deck rigs & wind tunnel and CFD of 2 to 5 units \citep{bordogna2020,mboumba2022,guzelbulut2024arrangement} & $-7$ to $+8$ points of thrust kept \\
 & Kite glide ratio & a soft kite in flight \citep{oehler2019} & $-3$\,\% \\
 & Kite flying height & SkySails' ship kites \citep{noaa2009} & inside their 100 to \SI{300}{m} \\
 & VTOL kite lowest wind & Makani's Wing~7 power curve \citep{vanderlind2013} & $+1.2$\,m/s \\
\addlinespace
Mass & Rotor mass per area & 5 Norsepower rotors \citep{norsepower2025}, which set it & $-12$ to $+3$\,\% \\
 & Wing sail mass per area & 3 wing sails built \citep{oceanwings2026web,riviera2024windwing,oceanbird2026wing560} & $-30$ to $+102$\,\% \\
 & Suction sail mass per area & 3 suction sails built \citep{conoship2022econowind,econowind2026vf3,econowind2026vf5} & 0 to $+116$\,\% \\
 & Wing sail foundation & the GT Wings sail of the Vectis Progress \citep{fednav2025gtwings} & about $-20$\,\% \\
 & Soft kite mass & SkySails kites of 160 and \SI{320}{kN} \citep{paulig2013} & $+4$ and $-4$\,\% \\
 & LEI kite mass & the TU Delft V2 kite \citep{schwoll2012}, whose tubes set it & $-37$\,\% \\
 & VTOL kite mass over size & wings built and designed \citep{sommerfeld2022scaling,vanhagen2023,noth2008} & exponent 2.4 to 2.6, inside their 2.2 to 3.1 \\
\addlinespace
Cost & Rotor price & 4 prices paid or quoted (\cref{tab:rotor_prices}) & $-15$ to $+3$\,\% \\
 & Canopy price & a sport kite \citep{varuste2026soul21} and army parafoils bought \citep{army2014ma7804} & above the kite, below the parafoils with their guidance \\
 & Shares of fuel, capital and crew in the ship's cost & published cost studies \citep{karvonen2020,healy2019} & inside their ranges \\
\bottomrule
\end{tabular}}
\end{table*}

\subsection{Whole ships at sea}
\label{sec:measured}

A ship measured at sea with its rig tests the whole model at once: the hull in calm water and waves, the wind on the hull, the drift and the rudder, the propeller and the engine at part load, and the rig's force. Few rigs have been measured at sea, and none on our routes. We model the strongest published measurement of each rig as closely as the model allows (\cref{tab:measured}). Ship, route, climate and the way each saving was measured differ from case to case, so a gap says something about the model only where these differences cannot account for it.

\paragraph{Rotor on a speed trial} A speed trial leaves route and climate out. The ferry Copenhagen, \given{\SI{156}{m}} between perpendiculars, carries one rotor \given{\SI{30}{m}} tall and \given{\SI{5}{m}} across and was sailed in four pairs of runs with the rotor turning and stopped \citep{werner2021speedtrial}. From the changes in speed and shaft power the trial report derives the rotor's thrust coefficient over the apparent wind angle, and the net power saving at \given{\SI{16}{kn}} in a true wind of \given{\SI{10}{m/s}} at \given{\SI{10}{m}}, which it takes as \given{\SI{11.8}{m/s}} at the rotor's mid-height \citep{werner2023copenhagen}. In the same apparent wind and spinning at the trial's \given{170} revolutions a minute, our rotor's thrust lies within \NumCopenhagenThrustBow\,\% of the measured with the apparent wind forward of \given{\SI{50}{\degree}}, but is \NumCopenhagenThrustBeam\,\% higher with the true wind just abaft the beam and \NumCopenhagenThrustAft\,\% higher with it on the quarter (\cref{tab:measured}a). The last point is the least sure: there the trial's rotor turned at a spin ratio of about \given{6}, beyond the \given{4} at which our coefficients stop, and the measured coefficient is the report's quadratic fit carried a little past its last trial point. With our thrust and spin power in place of the measured, the net saving stays as measured with the wind forward and rises by several points with the wind abeam and aft. The trial report finds the same pattern against a rotor computed with CFD on a plane with no ship under it, and puts it down to the hull and superstructure, which turn and slow the wind the rotor meets, most of all near the beam \citep{werner2023copenhagen}. Our rotor stands in the free wind (\cref{sec:deck_rigs}), so it carries this error: on this ferry it overstates the thrust with the wind abeam and aft by a quarter or more. How far the loss differs behind the lower deckhouse of a bulk carrier or tanker the trial does not say.

\paragraph{Rotors and wing sails in service} Four cargo ships report a fuel saving over months of service (\cref{tab:measured}b). TR Lady, an \given{\num{82000}}~t Kamsarmax bulk carrier with three rotors \given{\SI{24}{m}} by \given{\SI{5}{m}}, saved \given{9.1}\,\% of its propulsion fuel, \given{1.9}~t a day, over eight laden and ballast legs across the world's oceans, as computed by a performance model calibrated on more than \given{167} tests switching the rotors on and off and verified by Lloyd's Register \citep{lr2025trlady,lr2025onoff}. Maersk Pelican, a \given{\num{110000}}~t product tanker with two rotors \given{\SI{30}{m}} by \given{\SI{5}{m}}, saved \given{8.2}\,\% of its fuel over a year against a baseline from full-scale measurements and computation, which Lloyd's Register validated \citep{norsepower2019pelican}. Pyxis Ocean, a Kamsarmax with two rigid wing sails \given{\SI{37.5}{m}} tall, is reported to save \given{3}~t a day, \given{14}\,\% \citep{gcaptain2024pyxis}. Another report of the same results gives that figure as the prediction made before the trial and says the first six months came within \given{10}\,\% of it \citep{maritimeexec2024pyxis}, and the classification society's interim check covered sailing in favourable conditions only \citep{riviera2024windwing}. Shofu Maru, a \given{\num{100000}}~t coal carrier with one rigid sail up to \given{\SI{53}{m}} tall and \given{\SI{15}{m}} wide, saved \given{5} to \given{8}\,\% a voyage on voyages to Japan from Australia, Indonesia and North America, by a method Lloyd's Register verified \citep{mol2024shofu,mol2022shofu}. We sail the nearest class of ours on VLSFO with the same rig area against the same ship without a rig, at the class's operating speed on a fixed course, the shortest of a route of ours, as these ships held their speed rather than being routed: the bulk carriers to Tubarão and the tanker to Houston, over the \given{12} departures of the results. Shofu Maru, \given{\SI{235}{m}} long and \given{\SI{43}{m}} wide, takes the Aframax hull, the nearest in size, on the route to Tubarão. The model lays out each area by its own rule, so TR Lady's \given{\SI{360}{m^2}} become two rotors \given{\SI{33}{m}} tall and Shofu Maru's sail one wing \given{\SI{40}{m}} tall. Our fuel includes the hotel load, which TR Lady's figure leaves out.

The rotors save \NumMeasuredRotorGap{} points more of the fuel in the model than measured. In tonnes a day the model matches TR Lady: our Kamsarmax burns \NumMeasuredBareKamsarmax{}~t a day without a rig, against about \given{21}~t that TR Lady's two figures imply, so a like amount saved is a larger share of less. The rotors' gap points the way the Copenhagen trial does: rotors drive most with the wind abeam and aft, where our free-standing rotor is too strong. TR Lady and Maersk Pelican sailed from the tropics to the Arctic and the Southern Ocean, our ships the North Atlantic, and the climate may move the gap either way. The wing sails save \NumMeasuredPyxisGap{} points less than the figure reported for Pyxis Ocean and also come out below Shofu Maru's range. That figure is not yet an independently verified average over voyages, and unlike the rotor's the wing's thrust has no speed trial to check it against, so this gap is unexplained.

\paragraph{Kite} The one kite measured on a voyage is the \given{\SI{160}{m^2}} kite of the \given{\SI{132}{m}} Beluga SkySails on its first Atlantic crossing, from Bremen to Guanta in Venezuela between \given{22 January and 5 February 2008}. A crew still learning the system launched it almost daily and flew it \given{3} to \given{4} hours a day on average, and it saved half a tonne of the ship's \given{20}~t of heavy fuel oil a day, \given{2.5}\,\% \citep{verkehrsrundschau2008beluga}. We sail our \given{\SI{130}{m}} VLSFO ship at its service speed of \given{\SI{10.5}{kn}} on the shortest course out to Guadeloupe from every January departure, with the kite's nominal size read once as flat and once as projected area (\cref{tab:measured}b). It saves \NumBelugaTonnes~t a day and flies \NumBelugaFlightShare\,\% of the time. The fuel saved per day depends little on the hull, and there the model lies a little below the measurement; a kite's pull saves power in proportion to the ship's speed, so at the Beluga's \given{\SI{13}{kn}} the same pull would save about a quarter more. The model flies the kite far longer for a smaller gain per hour, since it flies wherever that pays, in light winds as well. Our ship burns less than half the Beluga's fuel, so the same tonnes are a larger share of it. The figures usually quoted for this ship are not measured. The project's record reports a saving of \given{5}\,\% on the ship's average route mix and \given{10 to 12}\,\% on North Atlantic and North Pacific routes \citep{wintecc2009_record}. Its public report rests these on \given{198}~t a year that a ship with the \given{\SI{160}{m^2}} kite can save on windy routes, and doubles that for the \given{\SI{320}{m^2}} kite, which was fitted only as the project ended \citep{wintecc2009}. SkySails computed such savings from the wind at about \given{\SI{200}{m}} along a ship's logged positions, with the kite flying on \given{30 to 50}\,\% of the days at sea \citep{skysails2011performance}: they say what the kite could have saved, not what it did. The one measured crossing gives no sign that the model favours the kite or holds it back.

If the gaps of the rotors and wing sails hold, the results favour the rotors over the wing sails by a few points of fuel saving.

\begin{table}[tbp]
\centering
\footnotesize
\caption{Our model beside rigs measured at sea. (a) The rotor of the ferry Copenhagen at \given{\SI{16}{kn}} in a true wind of \given{\SI{11.8}{m/s}} at its mid-height, from four true wind angles, spinning as on its speed trial: the apparent wind angle and spin ratio, the thrust coefficient on height times diameter as measured \citep{werner2023copenhagen} and of our rotor, and the net power saving as measured and with our rotor in its place. (b) Ships whose fuel saving was measured at sea, each as the nearest ship of ours with the same rig area on VLSFO at fixed speed on a fixed course: the rotor and wing-sail ships at the class's operating speed over the departures of the results, the Beluga SkySails' kite, read as flat and as projected area, at \given{\SI{10.5}{kn}} out to Guadeloupe from each January departure. The saving of fuel and the fuel saved per day at sea as reported, and ours with its standard error over the departures. Method study, \studyrun{measured_savings}.}
\label{tab:measured}
\adjustbox{max width=\columnwidth}{\begin{tabular}{@{}l@{}}\InputIfFileExists{sections/measured_table}{}{\tbd{Measured table pending: run the measured\_table step.}}\end{tabular}}
\end{table}

\subsection{Wind and waves}
\label{sec:val_weather}

The ships sail in ERA5's wind on a \given{\SI{1.5}{\degree}} grid every \given{6}~hours and in its waves on a \given{\SI{0.5}{\degree}} grid every \given{12}~hours, interpolated in between (\cref{sec:routing}). \Cref{tab:era5_buoys} sets them beside every record over \given{2017 to 2019} of the six moored buoys of the US National Data Buoy Center on or near our routes to the Caribbean, New York and Houston \citep{ndbc2026stdmet}, each buoy's wind carried from its anemometer, about \given{\SI{4}{m}} up, to \given{\SI{10}{m}} by the model's wind profile. ERA5's mean wind is \NumEraWindBias~m/s low, the mean of its square, which the rigs' forces follow, \NumEraWindSquare\,\% low, and winds from \given{\SI{12}{m/s}} \NumEraStrongWind\,\% low, since a coarse grid and six-hourly steps smooth out the strongest winds. Carried up by the model's steepest profile, the buoys' winds would be \given{4}\,\% stronger again. The mean wave height matches within \NumEraWaveBias~m, and seas from \given{\SI{3}{m}} are \NumEraHighWaves\,\% low. Weaker wind costs every rig thrust, so the savings are if anything low, the deck rigs' a little more than the kite's, which in strong winds is held at its force limit; seas slightly too low let the kite launch a little more often than it should. None of the buoys lies on the European end of the routes.

\begin{table*}[tbp]
\centering
\footnotesize
\caption{ERA5 as the routes sail in it beside every record of six moored buoys over \given{2017 to 2019} \citep{ndbc2026stdmet}: the buoy's mean \given{10\,m} wind, carried from its anemometer by the model's wind profile, and ERA5's bias against it, its mean square over the buoy's less one and its mean less the buoy's where the buoy measures \given{\SI{12}{m/s}} or more; the buoy's mean significant wave height, ERA5's bias and correlation with it, and its mean less the buoy's where the buoy measures \given{\SI{3}{m}} or more.}
\label{tab:era5_buoys}
\InputIfFileExists{sections/era5_buoy_table}{}{\tbd{ERA5 buoy table pending: run the era5\_buoy\_table step.}}
\end{table*}

\subsection{Ship}
\label{sec:val_ship}

\paragraph{Calm water} The worked example of \citet{holtrop1982} comes out within \given{0.2}\,\% in wetted surface, friction and wave-resistance coefficients, and its total resistance \given{1.2}\,\% higher, since we use the later form factor and wave term of \citet{holtrop1984}. Against towing-tank tests at model scale, without the correlation allowance that only the full-size ship carries, the method lies \given{0.3}\,\% above the Japan Bulk Carrier at its design speed \citep{t2015}, \given{2 to 5}\,\% above a Panamax tanker between \given{10} and \given{\SI{16}{kn}} at full scale and \given{10}\,\% above at \given{\SI{17}{kn}} \citep{farkas2017}, and \given{2 to 3}\,\% above the KCS container ship between \given{10} and \given{\SI{14}{kn}} and \given{6 to 10}\,\% from \given{18} to its service speed of \given{\SI{24}{kn}} \citep{t2015}. Its resistance thus grows \given{2 to 7}\,\% too fast between \given{\SI{10}{kn}} and service speed. Since each class is calibrated at its service speed, the calibrated hulls ask up to that much too little fuel when slow, which makes slowing down look cheaper than it is, for the ship with and without a rig alike. Beyond a prismatic coefficient of \given{0.85}, the limit tabulated for tankers and bulk carriers \citep{ortigosa2009}, the fitted wave term grows steeply; only the Kamsarmax lies beyond it, at \given{0.86}, so we evaluate its formulas at \given{0.85} with its true displacement; evaluated unchanged they would ask \given{3}\,\% less calm-water power at operating speed. The tank test of the Japan Bulk Carrier, at \given{0.84}, supports the method up to there.

\paragraph{Power of the classes} Uncalibrated, the model needs \given{4} to \given{14}\,\% more power than the source designs of the container ships, the car carrier, the Handysize and the MR and Aframax tankers, \given{6}\,\% less for the VLCC, and \given{19}, \given{27} and \given{32}\,\% more for the Capesize, Kamsarmax and Ultramax bulk carriers. One resistance factor per class, \given{0.80} to \given{1.05}, removes these gaps at service speed (\cref{app:calm_calibration}). Where the bulk carriers' excess lies, in their source powers, which MAN predicts for statistical average hulls \citep{man2022bulk}, or in the hull forms we estimate for them, is not known; were it all in the friction and form drag, their resistance would be up to \given{5}\,\% off one factor away from service speed. The coaster family keeps no factor, since its source power is a fleet average: its \given{\SI{85}{m}} coaster needs about \given{\SI{830}{kW}} for \given{\SI{10}{kn}}, where the Conoship CIP3800 of about its size makes that speed on \given{\SI{750}{kW}} by its builder's datasheet and the Damen Combi Freighter \given{3850} has a \given{\SI{1104}{kW}} main engine for \given{\SI{11.8}{kn}} \citep{maritimeexecutive2022}, so its resistance is if anything high.

\paragraph{Engine at part load} The engines' efficiency falls with the square of the load below full load, by the part-load factor $k=\given{0.2}$ fitted to MAN's L21/\given{31} medium-speed engine \citep[p.~137]{man2024l2131}. Against its four test-bed curves the model burns \given{1} to \given{7}\,\% more fuel per kilowatt-hour than the curves between \given{50 and 85}\,\% load, relative to full load, and \given{$-4$} to \given{$+6$}\,\% at \given{25}\,\%; a real engine burns least near \given{75 to 85}\,\% load, where the model's quadratic has no minimum. A rig unloads the engine, so the model takes a little too much of the fuel a rig saves at moderate load back again: this errs against every rig, the most against the one that unloads the engine most. The two-stroke engines of the large classes take the same law with their own peak efficiency; their part-load curves are not checked.

\paragraph{Drift and leeway}
\label{sec:val_drift}
\label{sec:drift_check}
Every wing sail and rotor result, and the kite's on courses across the wind, passes through the balance of hull, rudder and heel at the leeway the routed ships reach (\cref{sec:balance}). The kite is least exposed: it pulls mostly from behind the beam, where its side force is small. The link that matters most is the hull's side force. Captive tests of the tanker KVLCC2M at \given{3 to \SI{12}{\degree}} of leeway carry \given{0.55} times the slope of the regression of Kijima et al. at small leeway and more as the cross flow grows, which our factor $k_\beta$ and cubic term $q_\beta$ take from them, and place the side force about half a length forward of midship, moving aft as the leeway grows, which the model's centre follows within the spread of the repeated runs (\cref{tab:centre_check}; \citealp{toxopeus2011}). Towing-tank tests of two \given{\SI{138}{m}} hulls built without a skeg for wind propulsion carry only \given{0.32 and 0.37} times the regression, still less than the model gives them, and at the same side force pay about twice the model's induced drag (\cref{tab:drift_check}; \citealp{vanderkolk2021}). A skeg, bilge keels and the rudder raise the side force \citep{vanderkolk2021}, so the conventional single-screw hulls of our classes should lie between these hulls and the tanker, the only test of such a hull we found. With $k_\beta$ at the low bound of the hulls without a skeg, the kite loses \NumWeakHullKite, the wing sails \NumWeakHullWings{} and the rotors \NumWeakHullRotors{} points of cost saving on e-ammonia; this is the range the side-force results carry. The tests are at model scale, where the flow separates earlier than at full scale, and CFD of a cargo ship gives a drift-induced drag at full scale a third below that at scale \given{1:20} \citep{kramer2022b}, so the drag at the low bound covers the full-size hulls too. The rudder takes the lifting-line model of \citet{kramer2022b}, which gives the rudder drag of their CFD within \given{10}\,\% sailing straight and \given{25}\,\% on average; the model takes no effect of heel on the hull, which CFD at \given{\SI{20}{\degree}}, beyond our limit of \given{\SI{8}{\degree}}, finds lowers the side force by \given{7 to 20}\,\% \citep{kramer2022b}.

Studies that compare a free side force with a priced one find that it takes \given{8 to 26}\,\% of the fuel saving: a tanker with rotors keeps \given{12, 18 and 20}\,\% of fuel saving instead of \given{14, 23 and 27}\,\% at \given{300, 600 and \SI{900}{m^2}} \citep{tillig2019}, and a cargo ship with sails \given{17 to 65}\,\% instead of \given{21 to 71}\,\% \citep{kramer2022}. In our model it takes \NumSideForceRotors\,\% of the rotors' cost saving on e-ammonia, inside that range, and \NumSideForceWings\,\% of the wing sails', inside it on the \given{\SI{140}{m}} ship and above it on the Neo-Panamax container ship. The kite's saving moves by \NumSideForceKite\,\%, no more than the router's scatter.

\begin{table}[tbp]
\centering
\caption{Side-force coefficient on $\tfrac12\rho_w V^2 L T$ and centre of the side force forward of midship of the tanker KVLCC2M, from captive tests at scale \given{1:64.4} \citep{toxopeus2011}, against the model's hull of the same length, beam, draught and block coefficient at the same leeway; the two runs at \given{\SI{3}{\degree}} differ as the tests' three decimals allow.}
\label{tab:centre_check}
\InputIfFileExists{sections/centre_check_table}{}{\tbd{Centre check table pending: run the centre\_check\_table step.}}
\end{table}

\begin{table}[tbp]
\centering
\caption{Side-force coefficient and induced-drag coefficient at that side force, both on $\tfrac12\rho_w V^2 L T$, of two hulls built without a skeg for wind propulsion: the fit to towing-tank tests at scale \given{1:50} \citep{vanderkolk2021} against the model's hull of the same length, beam, draught and block coefficient.}
\label{tab:drift_check}
\InputIfFileExists{sections/drift_check_table}{}{\tbd{Drift check table pending: run the drift\_check\_table step.}}
\end{table}

\subsection{Rigs}
\label{sec:val_rigs}

The rotor's thrust is checked on the whole ship above. The wing sail's highest lift of \given{2.0}, for a wing that also trims its flap, lies above the \given{1.8} and \given{1.9} that a two-element wing with its flap held reached in the wind tunnel and in CFD \citep{wright2024wingsail}; no full-scale measurement of a wing sail's force is published. The suction sail's force coefficient is set from strain gauges on two VentiFoils at sea \citep{ghorbani2024suction}, and the kite's flight drag from the glide ratio of a soft kite in flight \citep{oehler2019}, so neither is checked by them.

\paragraph{Losses between deck rigs}
\label{app:interference}
Where units shade each other depends on the line joining them: units in line on the deck do so in head and following winds, units abreast in beam winds. In the wind tunnel, two rotors three diameters apart lost \given{4} to \given{6}\,\% of driving force on average over the wind angles of a reaching ship in line and \given{7} to \given{9}\,\% side by side, though up to \given{37}\,\% at \given{\SI{30}{\degree}} apparent wind in line \citep{bordogna2020}. CFD of wing sails on a bulk carrier gives \given{98}\,\% of the isolated thrust for three sails \given{2.4} chords apart and \given{77}\,\% for five sails \given{1.4} chords apart \citep{mboumba2022}, the chord being the depth of a wing from front to back; our array factor gives \given{91} and \given{85}\,\%, \given{7} points less and \given{8} more. Two-dimensional potential flow, which sets the model's factor for two columns, keeps \given{72} to \given{89}\,\% of the thrust for four wing sails in two columns \citep{vanreen2026}.

\subsection{Mass and cost}
\label{sec:val_mass}
\label{sec:val_cost}

\begin{figure}[tbp]
\centering
\plot{deck_mass}
\caption{Mass per area of wing sails, rotors and suction sails against their height: the size law of \cref{eq:size_law} at base inputs between the bounds of its exponent $k$, solid up to the tallest unit built and dashed beyond it, and the built units, rotors with their foundations \citep{norsepower2025}, wing sails \citep{oceanwings2026web,riviera2024windwing,oceanbird2026wing560} and suction sails \citep{conoship2022econowind,econowind2026vf3,econowind2026vf5}. A rotor's area is its height times its diameter, a sail's its height times its chord.}
\label{fig:deck_mass}
\end{figure}

\paragraph{Mass} Each deck rig weighs the mass per square metre of its built units up to the tallest built, and grows by the size law beyond (\cref{fig:deck_mass}). The model, set at the largest Norsepower rotor, gives the smallest \given{12}\,\% too little, since a drive, bearings and end disc of fixed mass weigh much of a small unit, and the others within \given{4}\,\% \citep{norsepower2025}. The wing sails scatter by a factor of three between makers, from \given{\SI{0.12}{t/m^2}} for the OceanWings \citep{oceanwings2026web} to \given{0.36} for the Oceanbird wing with its foundation \citep{oceanbird2026wing560}, and the model's \given{0.25} lies between them. The mass sets both the price and the cargo a rig displaces, so the structure of a heavy wing such as the Oceanbird would cost almost half as much again as we count, and that of a light one such as the OceanWings half. The soft kite's mass law gives SkySails' \given{160} and \given{\SI{320}{kN}} kites \given{4}\,\% more and \given{4}\,\% less than they weighed with their lines \citep{paulig2013}, the only kites built for a ship's force per area (\cref{fig:kite_mass}).

\paragraph{Rotor prices} Wing sails and kites have no public price we can use; only rotors do. \Cref{tab:rotor_prices} sets the rotor prices paid or quoted against our build-up price of the same rotors. Norsepower's 2019 figures are per installed unit \citep[p.~15]{tillig2020b}, so they are set against our build-up with its \given{15}\,\% for installation; so is Anemoi's, which sold its rotors supplied and installed \citep{maritimeexec2026anemoi}, mostly folding, so with the premium of a retractable rig as well, and is the mean over its \given{29} rotors of their sales value. Prices are brought to the euros of 2025 and to dollars at USD~\given{1.16} a euro. Our build-up lies \given{3}\,\% above the price of the smallest rotor and \given{9} to \given{15}\,\% under the other three, so the rotors' cost per area errs on balance towards a cheaper rotor, which favours the rotors slightly. Not used are the asset costs that \citet{emsa2023} estimated for each rig from interviews, a maker's statement of what its rig will cost, a price reported without a buyer or a source, a project budget holding the development of a first unit, and a partner's budget in a funded project; of the figures in public sources these are most, among them every one for a wing sail tied to a product \citep{positivenews2023wind,pbcf2020windmodern}.

\begin{table}[tbp]
\centering
\footnotesize
\caption{Every price paid or quoted for a rotor that we have read in full: its kind and year, the rotor's height and diameter, its price in thousands of today's USD by the euro area consumer price index \citep{eurostat2026hicp}, our build-up price of the same rotor, the price over the build-up, and the source of the price.}
\label{tab:rotor_prices}
\InputIfFileExists{sections/rotor_prices_table}{}{\tbd{Rotor prices table pending: run the rotor\_prices\_table step.}}
\end{table}

\paragraph{The ship's cost} Published cost splits leave out the interest on the cargo, which the cargo owner bears, and mostly the port dues. \stale{designs}{On that basis fuel is 30\,\% of the cost of our \SI{85}{m} ship and 41\,\% of our \SI{200}{m} ship on VLSFO, at the speeds chosen by cost of about 10 to 11~kn.} \Citet{karvonen2020} give \given{23}\,\% for a general cargo ship of \given{\num{3500}}~t deadweight at \given{\SI{9}{kn}} and \given{51}\,\% for a bulk carrier of \given{\num{60000}}~t at \given{\SI{11.6}{kn}}, with fuel about a sixth cheaper than ours. The day rates of \citet{healy2019} give \given{30} to \given{45}\,\% for a Handysize bulk carrier between \given{10} and \given{\SI{13}{kn}}, once their earnings count as the return on capital. Fuel's share grows with speed, so ours sits where these put it. \stale{designs}{Capital with maintenance and insurance is 48 to 50\,\% of our cost against 34 to 58\,\% in \citet{karvonen2020}, and crew and other running costs 9 to 17\,\% against their 11 to 13\,\%, larger on the smaller ship.} The share of fuel decides how much of a ship's cost a rig can save, so these agreeing means the savings are not inflated by a cost too heavy in fuel.

\subsection{Not yet checked}
\label{sec:val_open}

Some parts of the model have published measurements we have not yet set it beside, and some have none. The wind loads on our classes' hulls come from the wind-tunnel coefficients of their types \citep{blendermann1993}, about which measured single ships scatter, and the bulk carriers borrow the laden tanker's shape without their cranes and hatch coamings, the raised rims of the cargo hatches; no measurement of our classes' own windage tells which way they lie. The forecast is scored against ERA5 (\cref{fig:forecast}), not against observed weather. The two-stroke engines' part-load fuel, the heel and trim of a ship under a rig at sea, and the kite's winch, launch mast and control unit against the prices of comparable machines have published counterparts that a later check could use. Wing sails and suction sails have no public price tied to a product, and no ship kite a measured canopy life or loss rate, so these stay goals and estimates (\cref{tab:goals}).
